%% Formule de l'operateur : Combine_flag_EC_Foken
%% Dans ce fichier, les �ventuels accents sont cod�s
%% en ISO-8859-15 (aka latin9)
%% La ligne suivante donne trois exemples 
%% e accent aigu : �, e accent grave : �, e-dans-l'o : �

$Flag_{ITCmax}$ est la valeur maximale entre $Flag_{ITCX}$ et $Flag_{ITCY}$ \\
\\
La valeur du flag r�sultat est calcul� suivant le tableau ci-dessous : \\
\begin{center}
\begin{tabular}{lll}
\hline 
$Flag_{XY}$ & $Flag_{stat}$ & $ITC_{max}$ \\ 
\hline 
1 & 1 & 1-2 \\ 
\hline 
2 & 2 & 1-2 \\ 
\hline 
3 & 1-2 & 3-4 \\ 
\hline 
4 & 3-4 & 1-2 \\ 
\hline 
5 & 1-4 & 3-5 \\ 
\hline 
6 & 5 & $\leq$ 5 \\ 
\hline 
7 &  $\leq$ 6 &  $\leq$ 6 \\ 
\hline 
8 & $\leq$ 8 &  $\leq$ 8 \\ 
\hline 
9 & 9  &  - \\ 
\hline 
\end{tabular} 
\end{center}
\vspace{0.8cm}
Indications d'utilisation des mesures de la covariance $\overline{X'Y'}$ en fonction de la valeur de $Flag_{XY}$ :
\begin{itemize}
\item de 1 � 3  : les mesures peuvent �tre utilis�es pour des applications n�cessitant la meilleure pr�cision comme la recherche et le d�veloppement de param�trisations. 
\item de 1 � 6 : les mesures peuvent �tre plut�t utilis�es pour des �tudes long terme.
\item de 7 � 8 : les mesures doivent �tre utilis�es uniquement comme ordre de grandeur.
Elles peuvent �tre une alternative � un traitement de "gap filling" si l'�cart avec les valeurs avant et apr�s n'est pas trop important.  
\item �gal � 9 : mesures � rejeter
\end{itemize}
